% =====================================================================
% Mirante dos Dados — Working Paper #3
% Acesso desigual: cirurgia uroginecológica no SUS como indicador de
% pobreza estrutural e os limites da compensação fiscal por emendas
% parlamentares (2008–2025).
%
% Autor:  Leonardo Chalhoub (engenharia de dados, análise, redação)
% Versão: 2.0  —  Abril de 2026
% Compilação: pdflatex articles/uropro-serie-2008-2025.tex (2x)
%
% Este Working Paper tem como base a pesquisa original de TATIELI
% (2022, especialização em Enfermagem). Usa-se a cirurgia
% uroginecológica como INDICADOR LIMPO de oferta hospitalar
% especializada (procedimento eletivo, seguro, bem definido em
% SIGTAP) e cruza-se com agregados por UF do Programa Bolsa Família
% (Mirante WP n. 2) e da execução de Emendas Parlamentares (Mirante
% WP n. 1) para responder uma pergunta de federalismo: pobreza
% estrutural produz desigualdade de acesso, e o fluxo fiscal
% político compensa essa desigualdade?
% =====================================================================

\documentclass[12pt,a4paper,oneside]{article}

\usepackage[T1]{fontenc}
\usepackage[utf8]{inputenc}
\usepackage[brazilian]{babel}
\usepackage{newtxtext}
\usepackage{newtxmath}
\usepackage[section]{placeins}    % auto-FloatBarrier no início de cada \section
\usepackage[a4paper,top=3cm,bottom=2cm,left=3cm,right=2cm]{geometry}
\usepackage{setspace}
\usepackage{indentfirst}
\usepackage{titlesec}
\usepackage{caption}
\usepackage{booktabs}
\usepackage{array}
\usepackage{multirow}
\usepackage{graphicx}
\usepackage[hidelinks]{hyperref}
\usepackage{url}
\usepackage{float}
\usepackage{xcolor}

\onehalfspacing
\setlength{\parindent}{1.25cm}
\setlength{\parskip}{0pt}

\titleformat{\section}{\normalfont\bfseries\large\uppercase}{\thesection}{0.6em}{}
\titleformat{\subsection}{\normalfont\bfseries\normalsize}{\thesubsection}{0.6em}{}
\titlespacing*{\section}{0pt}{18pt}{8pt}
\titlespacing*{\subsection}{0pt}{12pt}{6pt}

\captionsetup[table]{labelfont=bf,name=Tabela,labelsep=endash,
                     position=above,justification=centering,font=small}
\captionsetup[figure]{labelfont=bf,name=Figura,labelsep=endash,
                      position=below,justification=centering,font=small}
\captionsetup{skip=6pt}    % respiro fig/tab → caption (padrão Lancet/Cad SP)

\newcommand{\rs}{R\$\,}

% Stamp de compilação injetado pelo CI (deploy-pages.yml).
% Para builds locais, este arquivo default é usado — basta um pdflatex direto
% e a capa exibe "(dev local)". Em CI, este arquivo é sobrescrito antes da
% compilação com timestamp BRT + SHA do commit que disparou o deploy.
\newcommand{\COMPILESTAMP}{compila\c{c}\~ao local}
\newcommand{\COMPILESHA}{dev}


\begin{document}

% ---------- CAPA ---------------------------------------------------------
\begin{titlepage}
  \begin{center}
    \vspace*{1.5cm}
    {\small\textsc{Mirante dos Dados \quad\textbar\quad Working Paper n. 3}}\\
    {\small Versão 2.0 \textemdash{} Abril de 2026}\\[2pt]
    {\footnotesize Compilada em \COMPILESTAMP\ \textbullet\ \texttt{\COMPILESHA}}

    \vfill

    {\LARGE\bfseries
     ACESSO DESIGUAL:
     CIRURGIA UROGINECOLÓGICA NO SUS
     COMO INDICADOR DE POBREZA ESTRUTURAL
     E OS LIMITES DA COMPENSAÇÃO FISCAL
     POR EMENDAS PARLAMENTARES
     (2008\textendash 2025)}

    \vspace{1cm}
    {\large\itshape
     Unequal access: uro-gynecological surgery in Brazil's Unified
     Health System as an indicator of structural poverty and the
     limits of political-fiscal compensation through parliamentary
     amendments, 2008\textendash 2025}

    \vfill

    {\large\textbf{Leonardo Chalhoub}\textsuperscript{1}}\\[6pt]
    {\small\textsuperscript{1}\,Engenharia de Dados \textemdash{}
     análise quantitativa, integração entre verticais, redação.}\\[6pt]
    {\small\href{mailto:leonardochalhoub@gmail.com}{leonardochalhoub@gmail.com}
     \quad\textbullet\quad
     \href{https://github.com/leonardochalhoub}{github.com/leonardochalhoub}}

    \vfill

    \rule{0.6\linewidth}{0.4pt}\\[6pt]
    {\small\textit{Este artigo é o terceiro Working Paper do Mirante dos
     Dados. Diferentemente da maioria dos papers da série, que
     analisa uma vertical isoladamente, este é um trabalho
     deliberadamente \textbf{transversal}: usa a vertical UroPro
     (cirurgia uroginecológica) como indicador clínico de oferta
     hospitalar e cruza com as verticais Bolsa Família (cobertura
     PBF como proxy de pobreza estrutural) e Emendas Parlamentares
     (execução por UF como proxy de fluxo fiscal político) para
     interrogar a relação federalista entre pobreza, fluxo
     compensatório e acesso a serviços públicos especializados.
     Tem como base a pesquisa original de TATIELI (2022,
     especialização em Enfermagem).}}

    \vspace{1cm}
    {\footnotesize
     Disponível em
     \href{https://leonardochalhoub.github.io/mirante-dos-dados-br/}%
     {leonardochalhoub.github.io/mirante-dos-dados-br}}\\[2pt]
    {\footnotesize
     \href{https://github.com/leonardochalhoub/mirante-dos-dados-br}%
     {github.com/leonardochalhoub/mirante-dos-dados-br}}
  \end{center}
\end{titlepage}

% =====================================================================
\section*{Resumo}
\addcontentsline{toc}{section}{Resumo}
% =====================================================================
\begin{singlespace}\small\noindent
\textbf{Objetivo.} Investigar se a desigualdade territorial de
acesso à cirurgia uroginecológica no SUS reflete pobreza
estrutural por UF e se o fluxo fiscal compensatório
\textemdash{} representado pela execução de emendas parlamentares
\textit{per capita} \textemdash{} corrige, atenua ou agrava esse
gradiente. A cirurgia para tratamento da incontinência urinária
é usada como \textit{indicador clínico} pela conjunção de três
propriedades: é eletiva (sensível à oferta hospitalar), é
clinicamente segura (mortalidade abaixo de 0{,}05\,\%, removendo
o risco como variável de seleção) e tem definição operacional
estável via SIGTAP.

\textbf{Métodos.} Dados de três verticais do Mirante dos Dados:
(i) UroPro \textemdash{} microdados SIH-AIH-RD/DATASUS,
procedimentos SIGTAP 0409010499 (via abdominal) e 0409070270
(via vaginal), 2008\textendash 2025; (ii) Bolsa Família \textemdash{}
agregados estaduais de beneficiários e valores transferidos,
2013\textendash 2025; (iii) Emendas Parlamentares \textemdash{}
execução orçamentária por UF (RP6, RP7, RP9), 2014\textendash 2025.
Despesas deflacionadas por IPCA (BCB) para \rs 2021. Populações
do IBGE/SIDRA. As três fontes passam pela mesma arquitetura
medalhão (\textit{bronze}\,/\,\textit{silver}\,/\,\textit{gold})
do Mirante; o \textit{gold} de cada vertical é exportado em
JSON versionado no repositório aberto.

\textbf{Resultados.} (i)~Em 2025, a taxa de cirurgias
uroginecológicas por 100\,mil habitantes varia de 14{,}71 (SC)
a 0{,}14 (RR), uma diferença de aproximadamente 100 vezes.
(ii)~A correlação de Pearson entre a cobertura PBF
(\%~da população beneficiária) e o acesso à cirurgia (AIH/100k)
é de aproximadamente \textendash 0{,}68 sobre as 27 UFs:
estados mais pobres têm sistematicamente menor acesso.
(iii)~A execução de emendas parlamentares \textit{per capita}
correlaciona-se \textit{negativamente} com o acesso à cirurgia
($\rho \approx -0{,}45$): UFs mais pobres recebem mais emendas
\textit{per capita}, mas isso não se traduz em ampliação visível
de oferta cirúrgica especializada. (iv)~Três UFs operam acima
do esperado dado seu nível de pobreza (TO, AC, MA), sugerindo
que existem configurações organizacionais locais capazes de
\textit{driblar} o gradiente. (v)~A pandemia de COVID-19
amplificou temporariamente, mas não criou, a desigualdade.
(vi)~A permanência hospitalar caiu 40\,\% na via vaginal entre
2008 e 2025 (eficiência clínica), evidência de que a barreira
ao acesso não é técnica.

\textbf{Conclusões.} A desigualdade de acesso à cirurgia
uroginecológica no SUS é \textit{estrutural} \textendash{}
acompanha o gradiente de pobreza estadual \textendash{} e
\textit{não é compensada} pelo mecanismo de emendas
parlamentares, apesar de seu volume nominal expressivo em UFs
periféricas. O federalismo brasileiro produz, neste caso,
universalismo formal e desigualdade material. A hipótese de
que a transferência fiscal política seja capaz de gerar oferta
hospitalar especializada deve ser submetida a evidência empírica
mais cuidadosa, não tratada como axioma.

\vspace{4pt}
\textbf{Palavras-chave:} SUS; incontinência urinária; cirurgia
eletiva; equidade em saúde; Bolsa Família; emendas parlamentares;
federalismo brasileiro; arquitetura medalhão; análise
\textit{cross-vertical}.
\end{singlespace}

% =====================================================================
\section*{Abstract}
\addcontentsline{toc}{section}{Abstract}
% =====================================================================
\begin{singlespace}\small\noindent
\textbf{Objective.} To investigate whether territorial inequality
in access to uro-gynecological surgery within Brazil's Unified
Health System (SUS) reflects state-level structural poverty,
and whether compensatory fiscal flow \textemdash{} represented
by per-capita execution of parliamentary amendments \textemdash{}
corrects, attenuates or worsens this gradient. Surgery for
urinary incontinence is used as a \textit{clinical indicator}
because it is elective (sensitive to hospital supply), clinically
safe (in-hospital mortality below 0.05\,\%, removing risk as a
selection variable) and operationally well-defined under SIGTAP.

\textbf{Methods.} Data from three Mirante dos Dados verticals:
(i) UroPro \textemdash{} SIH-AIH-RD/DATASUS microdata,
2008\textendash 2025; (ii) Bolsa Família \textemdash{} per-state
beneficiary and transfer aggregates, 2013\textendash 2025;
(iii) Parliamentary Amendments \textemdash{} per-state budget
execution, 2014\textendash 2025. All deflated to constant 2021
reais (IPCA, Brazilian Central Bank); state populations from
IBGE/SIDRA.

\textbf{Results.} (i)~In 2025, the per-100k surgical rate spans
14.71 (SC) to 0.14 (RR), a roughly 100-fold gap.
(ii)~Pearson correlation between PBF coverage and surgical
access is $\rho \approx -0.68$ across 27 states: poorer states
have systematically less access. (iii)~Per-capita parliamentary
amendments correlate \textit{negatively} with surgical access
($\rho \approx -0.45$): poorer states receive more amendments
per capita, but this does not translate into visible expansion
of specialized surgical supply. (iv)~Three states (TO, AC, MA)
overperform their poverty profile, suggesting that local
organizational configurations can override the gradient. (v)~The
COVID-19 pandemic temporarily amplified, but did not create, the
inequality. (vi)~Hospital length of stay fell 40\,\% on the
vaginal approach between 2008 and 2025, indicating clinical
efficiency \textemdash{} the barrier to access is not technical.

\textbf{Conclusions.} Access inequality is structural,
poverty-aligned, and not compensated by parliamentary
amendments. Brazilian federalism here produces formal
universalism alongside material inequality. The implicit
assumption that political-fiscal transfer generates specialized
hospital supply should be subjected to empirical scrutiny, not
treated as axiom.

\vspace{4pt}
\textbf{Keywords:} SUS; urinary incontinence; elective surgery;
health equity; Bolsa Família; parliamentary amendments; Brazilian
federalism; medallion architecture; cross-vertical analysis.
\end{singlespace}

% =====================================================================
\section{Introdução}\label{sec:intro}
% =====================================================================

O Sistema Único de Saúde brasileiro foi desenhado, na Constituição
de 1988, sobre o princípio da \textit{universalidade}: a saúde é
direito de todos e dever do Estado. Em sua execução real, contudo,
o SUS opera dentro de um federalismo de implementação altamente
heterogêneo, em que a oferta de serviços de média e alta
complexidade depende de redes hospitalares estaduais e municipais
com capacidades muito diferentes entre si. Essa tensão entre
\textit{universalismo formal} (todos têm direito) e
\textit{disponibilidade material} (a oferta é desigualmente
distribuída) é, há décadas, um dos eixos centrais do debate sobre
política pública de saúde no Brasil.

A literatura epidemiológica documenta amplamente o gradiente
saúde\,/\,pobreza tanto entre regiões quanto entre estratos
socioeconômicos. O que é menos comum é a quantificação do
gradiente sobre um \textit{procedimento clínico específico, bem
definido, eletivo e seguro}. A maior parte das comparações
estaduais usa indicadores agregados (mortalidade infantil,
expectativa de vida, cobertura de atenção primária) que
combinam muitas dimensões de oferta, demanda e composição
demográfica. Quando se quer interrogar especificamente
\textit{capacidade hospitalar especializada}, faz sentido adotar
um indicador clínico mais limpo.

Este artigo propõe usar a \textbf{cirurgia uroginecológica para
tratamento de incontinência urinária} (procedimentos SIGTAP
0409010499, via abdominal, e 0409070270, via vaginal) como
indicador limpo de oferta hospitalar especializada. A escolha é
deliberada e se justifica por três propriedades:

\begin{enumerate}
\item \textbf{É eletiva.} Cirurgias eletivas, por definição,
podem ser adiadas. Se a oferta é insuficiente, o paciente espera;
se é suficiente, o paciente é operado. O volume observado é,
portanto, sensível diretamente à oferta hospitalar e não é
artificialmente sustentado por urgência clínica.

\item \textbf{É segura.} A literatura clínica internacional
(FORD \textit{et al.}, 2017; CHAPPLE \textit{et al.}, 2020)
e os próprios dados do SIH-RD mostram mortalidade
\textit{intra-hospitalar} estruturalmente abaixo de 0{,}05\,\%.
Isso elimina \textit{risco clínico} como variável de seleção:
diferenças regionais de acesso não são explicadas por diferentes
perfis de risco do paciente.

\item \textbf{É bem definida.} Os dois códigos SIGTAP capturam o
universo da prática cirúrgica brasileira para a condição. Não há
ambiguidade significativa de classificação, o que permite
comparações diretas entre UFs e ao longo do tempo.
\end{enumerate}

Cumpridas essas propriedades, a desigualdade observada de acesso
à cirurgia uroginecológica entre UFs pode ser lida quase
exclusivamente como desigualdade de oferta.

A pergunta-mestra deste Working Paper é, então: \textit{a oferta
de cirurgia uroginecológica especializada acompanha o gradiente
de pobreza dos estados brasileiros, e os mecanismos
político-fiscais de compensação \textemdash{} as emendas
parlamentares \textemdash{} corrigem esse gradiente?}

Aproveitamos uma conjunção operacional rara: o Mirante dos Dados
mantém, na mesma plataforma e com a mesma arquitetura de pipeline,
três verticais cujos agregados estaduais são diretamente
comparáveis: UroPro (cirurgias por UF$\times$ano), Bolsa Família
(beneficiários e valores por UF$\times$ano) e Emendas Parlamentares
(execução por UF$\times$ano). O cruzamento dessas três séries é
trivial em termos computacionais e foi possibilitado pela
unificação metodológica da plataforma. Esse Working Paper, em
parte, é a demonstração de que a integração de verticais
\textit{não é apenas} engenharia de dados \textit{a posteriori}:
ela cria possibilidades analíticas novas.

Antecipamos a estrutura do argumento. A Seção~\ref{sec:teoria}
apresenta o marco conceitual que organiza os três indicadores.
A Seção~\ref{sec:metodo} descreve os dados, o pipeline e
\textemdash{} de forma transparente \textemdash{} um defeito
específico de pipeline descoberto e corrigido em abril de 2026,
que afetou as visualizações da plataforma até essa correção. A
Seção~\ref{sec:resultados} apresenta as evidências em quatro
recortes: caracterização da oferta cirúrgica, gradiente de
pobreza, fluxo fiscal compensatório e o cruzamento entre os três.
A Seção~\ref{sec:discussao} discute as implicações para o debate
sobre desigualdade, federalismo e política pública. A
Seção~\ref{sec:conclusao} fecha.

% =====================================================================
\section{Marco Conceitual}\label{sec:teoria}
% =====================================================================

\subsection{Universalismo formal e federalismo material}

A Constituição brasileira de 1988 institui o SUS como sistema
universal: \textit{toda} pessoa tem, em princípio, direito ao
atendimento de saúde pelo Estado. A execução do princípio,
contudo, é materialmente delegada a uma rede tripartite
(União\,/\,estados\,/\,municípios) cuja capacidade instalada é
historicamente desigual. Estados do Sul e Sudeste, com maior
densidade hospitalar, maior tradição de formação médica e maior
arrecadação tributária, oferecem sistematicamente mais
\textit{capacidade efetiva} do que estados do Norte e Nordeste.
Essa assimetria não é segredo \textemdash{} é parte da literatura
clássica sobre federalismo brasileiro \textemdash{} mas é, em
geral, descrita por proxies agregados (PIB \textit{per capita},
densidade médica, leitos\,/\,1\,000 hab.) cuja interpretação
clínica direta é difícil.

\subsection{Cirurgia eletiva como indicador clínico de oferta}

Procedimentos eletivos têm uma propriedade analítica útil:
sua quantidade observada é uma função quase direta da oferta
hospitalar disponível. Se há equipe treinada, leito e materiais
disponíveis, o paciente é operado; caso contrário, fica em
fila ou é redirecionado para o setor privado. Não há demanda
inelástica que sustente o volume artificialmente. Isso contrasta
com procedimentos de urgência, cuja realização é praticamente
mandatória (cesarianas de emergência, traumatologia,
internações por infarto), e por isso menos sensíveis à
disponibilidade do parque hospitalar local.

Dentro da categoria eletiva, a cirurgia uroginecológica para IU
tem uma vantagem adicional: \textit{baixa mortalidade
intra-hospitalar} significa que diferenças de acesso não podem
ser explicadas por diferentes perfis de risco do paciente.
Ou seja, se SC opera 14{,}7 cirurgias por 100k e RR opera 0{,}14,
a diferença não é porque pacientes de SC têm menos risco
operatório \textemdash{} é simplesmente porque há mais oferta
em SC.

\subsection{O Bolsa Família como proxy de pobreza estrutural}

A cobertura do Programa Bolsa Família (\%~da população estadual
em situação de beneficiário) é um proxy razoável de
\textit{pobreza estrutural} por duas razões. Primeiro, o
critério de elegibilidade é nacional e relativamente estável
(renda \textit{per capita} familiar abaixo de uma linha
indexada), o que torna a comparação entre UFs e ao longo do
tempo mais consistente do que indicadores derivados de pesquisa
amostral. Segundo, o tempo médio de permanência no programa é
longo (vários anos), o que faz da cobertura um indicador de
pobreza \textit{cronificada}, não conjuntural. Reconhecemos que
o PBF não é uma medida perfeita de pobreza \textemdash{} ele
captura também heterogeneidade na capacidade local do CRAS de
inscrever famílias \textemdash{} mas como proxy comparativo
entre UFs ele é defensável e, mais importante, está disponível
de forma agregada e atualizada na plataforma Mirante (CHALHOUB,
2026b).

\subsection{Emendas parlamentares como proxy de fluxo fiscal
político}

A execução de emendas parlamentares \textit{per capita} captura
um aspecto específico do fluxo fiscal federativo: recursos
alocados não pelo Executivo, mas por iniciativa de
parlamentares (deputados federais e senadores) com
discricionariedade sobre destino e finalidade. Há quatro tipos
principais (RP6 individual, RP7 bancada, RP8 comissão, RP9
relator); todos são endógenos ao processo legislativo, e o
montante \textit{per capita} recebido por estado depende de
cálculos políticos que não são proporcionais ao tamanho da
população nem ao grau de privação. Para o argumento deste
Working Paper, é suficiente que as emendas \textit{per capita}
representem \textit{algum} fluxo de recurso fiscal alocado por
processo político, distinto do orçamento regular do Ministério
da Saúde. A pergunta é: esse fluxo compensatório se traduz em
oferta hospitalar especializada visível? Investigamos
empiricamente.

% =====================================================================
\section{Aspectos Metodológicos}\label{sec:metodo}
% =====================================================================

\subsection{Fontes de dados}

\textbf{(i)~UroPro.} Microdados SIH-AIH-RD do DATASUS,
distribuídos via servidor FTP em arquivos DBC mensais por UF,
janeiro/2008 a dezembro/2025. Filtrados por procedimentos
SIGTAP 0409010499 e 0409070270.

\textbf{(ii)~Bolsa Família.} Pagamentos por UF e ano, série
2013\textendash 2025, originária do Portal da Transparência
da CGU, processada conforme descrito em CHALHOUB (2026b,
Working Paper n. 2). Indicadores derivados: número de
beneficiários, valor nominal, valor 2021 (deflacionado),
valor por beneficiário, valor \textit{per capita}.

\textbf{(iii)~Emendas Parlamentares.} Execução orçamentária por
UF e ano, série 2014\textendash 2025, do Portal da
Transparência, conforme CHALHOUB (2026a, Working Paper n. 1).
Categorias incluem RP6 (individual), RP7 (bancada), RP9
(relator) e outras. Valores deflacionados para \rs 2021 via
IPCA.

\textbf{(iv)~Auxiliares.} IPCA (BCB-SGS\,433) para deflação;
estimativas populacionais IBGE\,/\,SIDRA tabela 6579.

\subsection{Pipeline de dados (arquitetura medalhão)}

As três verticais passam pela mesma estrutura
(\textit{bronze}\,/\,\textit{silver}\,/\,\textit{gold})
descrita em CHALHOUB (2026a, 2026b). Em síntese:
\textit{bronze} preserva o dado bruto convertido para Delta
em formato \textit{string-only} (sem inferência de tipos);
\textit{silver} faz limpeza, normalização e agregação por
chave analítica; \textit{gold} entrega a vista colapsada por
(uf$\times$ano), enriquecida com população e deflatores. O
\textit{gold} é exportado em JSON versionado no repositório
público do projeto, permitindo reprodutibilidade externa.

\subsection{O defeito da pipeline silver e sua correção}\label{sub:bug}

Em abril de 2026, durante uma inspeção visual do mapa de
\textit{Distribuição geográfica} da plataforma na vertical
UroPro, observou-se que diversas unidades da federação
apareciam como \textit{sem dado} para combinações
(ano, procedimento) em que o sistema deveria, em princípio,
ter registros. Uma sequência de consultas SQL diretas ao Delta
da camada \textit{silver} revelou que, para o procedimento de
maior volume nacional (0409070270, via vaginal), apenas 13 das
27 UFs apareciam por ano \textemdash{} mesmo nos anos em que o
\textit{bronze} continha registros para todas as 27.

A causa raiz, identificada por inspeção do código, foi a
seguinte linha no \textit{notebook silver}:

\begin{singlespace}\small
\begin{verbatim}
latest_ts = bronze.agg(F.max("_ingest_ts")).first()[0]
src = bronze.where(F.col("_ingest_ts") == latest_ts)
\end{verbatim}
\end{singlespace}

\noindent
A intenção declarada (em comentário) era \textit{evitar dupla
contagem se rodaram o ingest mais de uma vez}. O bug está na
hipótese implícita: o filtro pressupõe que \textit{cada} valor
de \texttt{\_ingest\_ts} representa um \textit{snapshot}
completo do \textit{bronze}. Isso não é verdade. O Auto Loader
processa arquivos em micro-\textit{batches} dentro de uma mesma
execução, e \textit{cada} micro-batch recebe seu próprio
\texttt{\_ingest\_ts}. No caso real, os 408 arquivos por UF
foram distribuídos em quatro micro-batches separados por
poucos segundos. O filtro \texttt{== max} preservava
\textit{apenas} o último micro-batch, descartando 73\,\% das
linhas \textit{bronze} \textit{e} 14 das 27 UFs em silêncio.

A correção foi remover o filtro. A idempotência por arquivo é
garantida pelo \textit{checkpoint} do próprio Auto Loader, o
que torna o filtro \texttt{max(\_ingest\_ts)} não apenas
defeituoso mas também redundante. O \textit{commit} público
que aplica a correção é o \texttt{fa869cf} no repositório do
projeto. Toda a camada \textit{silver} e \textit{gold} foi
regenerada em seguida; os números apresentados neste Working
Paper já refletem essa reconstituição. \textbf{Magnitudes
absolutas anteriores à correção devem ser consideradas
prejudicadas, embora as direções de tendência se mantenham}
\textemdash{} o que é, por si, um achado metodológico relevante.

\subsection{Indicadores construídos}

Por (uf$\times$ano), trabalhamos com:

\begin{itemize}
\item \texttt{aih\_por\_100k}: cirurgias uroginecológicas (via
vaginal) por 100\,mil habitantes. Indicador-chave de oferta
hospitalar especializada.
\item \texttt{pct\_pop\_em\_pbf}: \%~da população estadual
beneficiária do PBF. Proxy de pobreza estrutural.
\item \texttt{emenda\_per\_capita\_2021}: execução de emendas
\textit{per capita}, em \rs 2021. Proxy de fluxo fiscal
político.
\item \texttt{dias\_perm\_avg}: permanência hospitalar média
ponderada por AIH. Indicador adicional de eficiência clínica.
\item \texttt{mortalidade}: óbito intra-hospitalar como fração
das AIH. Controle de risco clínico.
\end{itemize}

\subsection{Limitações}

(i)~A análise é \textit{descritiva e correlacional} \textemdash{}
correlação entre cobertura PBF e acesso à cirurgia não estabelece
causalidade, e em particular não exclui a possibilidade de uma
terceira variável omitida (densidade médica regional, por
exemplo) ser a causa comum de ambos. (ii)~A IU cirúrgica é
predominantemente um procedimento feminino; a comparação por UF
deveria, idealmente, ser ajustada por estrutura etária da
população feminina, o que não fazemos aqui. (iii)~O PBF mede
pobreza segundo um critério administrativo que mudou várias
vezes ao longo da série; usamos cobertura média 2024\textendash 2025
para minimizar essa heterogeneidade. (iv)~A execução de emendas
parlamentares por UF inclui recursos não destinados à saúde, o
que enfraquece a hipótese do mecanismo causal direto entre
emendas e oferta cirúrgica; tratamos as emendas como proxy de
fluxo fiscal político \textit{geral}, não específico do setor.

% =====================================================================
\section{Resultados}\label{sec:resultados}
% =====================================================================

\subsection{Caracterização da oferta cirúrgica
(2008\textendash 2025)}\label{sub:oferta}

A Tabela~\ref{tab:nat-aih} consolida o volume nacional anual de
cirurgias uroginecológicas, separado por via.

\begin{table}[H]
\centering
\caption{Volume nacional anual de cirurgias uroginecológicas por
via cirúrgica, 2008\textendash 2025. Fonte: SIH-AIH-RD/DATASUS,
processado pelo Mirante dos Dados (camada \textit{gold} corrigida).}
\label{tab:nat-aih}
\small
\begin{tabular}{rrrr}
\toprule
\textbf{Ano} & \textbf{Abdominal} & \textbf{Vaginal} & \textbf{Total} \\
\midrule
2008 & 2\,074 & 6\,558  & 8\,632 \\
2010 & 1\,542 & 6\,324  & 7\,866 \\
2012 & 1\,216 & 5\,982  & 7\,198 \\
2014 & 1\,032 & 5\,785  & 6\,817 \\
2016 &   789  & 5\,467  & 6\,256 \\
2018 &   760  & 5\,555  & 6\,315 \\
2019 &   759  & 5\,959  & 6\,718 \\
2020 &   384  & 2\,566  & 2\,950 \\
2021 &   323  & 2\,511  & 2\,834 \\
2022 &   646  & 5\,578  & 6\,224 \\
2023 &   821  & 6\,785  & 7\,606 \\
2024 &   863  & 8\,966  & 9\,829 \\
2025 &   740  & 10\,144 & 10\,884 \\
\bottomrule
\end{tabular}
\end{table}

Três fenômenos saltam à vista. Primeiro, a via vaginal predomina
de forma estável e crescente: a razão vaginal:abdominal sobe de
3{,}2:1 em 2008 para 13{,}7:1 em 2025, refletindo a substituição
clínica progressiva da via abdominal pela vaginal (\textit{slings}
sub-uretrais TVT/TOT) consagrada pela literatura
internacional. Segundo, o volume total apresenta queda gradual
entre 2008 e 2018 (de aproximadamente 8\,600 para 6\,300\,/\,ano)
seguida por recuperação e expansão até atingir, em 2025, o
patamar mais alto da série completa (10\,884). Terceiro, e mais
relevante para o argumento deste paper, é a queda de 56\,\% em
2020\textendash 2021 (relativamente à média
2015\textendash 2019, $\sim$\,6\,070\,/\,ano) seguida de
recuperação e \textit{ultrapassagem} do patamar pré-pandêmico
em 2024\textendash 2025. Esse pico é evidência de \textit{represa
cirúrgica} \textemdash{} estoque de fila acumulado durante a
pandemia em escoamento nos anos seguintes.

A queda da permanência hospitalar (Tabela~\ref{tab:nat-perm}) é
o achado clínico-organizacional mais relevante isoladamente: a
via vaginal vai de 2{,}39 dias\,/\,AIH em 2008 para 1{,}43 em
2025 \textemdash{} redução de 40{,}2\,\%, ou ganho de
produtividade hospitalar de 67\,\%. A via abdominal segue padrão
similar (\textendash 26{,}7\,\%). Esses ganhos são compatíveis
com a difusão de técnicas minimamente invasivas e indicam que a
\textit{eficiência} da prática hospitalar brasileira para esse
procedimento não está atrás da literatura internacional. A
barreira ao acesso, portanto, não é técnica.

\begin{table}[H]
\centering
\caption{Permanência hospitalar média nacional, em dias\,/\,AIH,
ponderada por número de internações. Fonte: SIH-AIH-RD.}
\label{tab:nat-perm}
\small
\begin{tabular}{rrr}
\toprule
\textbf{Ano} & \textbf{Abdominal (dias)} & \textbf{Vaginal (dias)} \\
\midrule
2008 & 2{,}47 & 2{,}39 \\
2012 & 2{,}22 & 2{,}22 \\
2015 & 2{,}24 & 1{,}92 \\
2019 & 2{,}14 & 1{,}86 \\
2021 & 1{,}94 & 1{,}55 \\
2023 & 1{,}79 & 1{,}53 \\
2025 & 1{,}81 & 1{,}43 \\
\midrule
\textbf{$\Delta$ 2008$\rightarrow$2025} & \textbf{\textendash 26{,}7\,\%} & \textbf{\textendash 40{,}2\,\%} \\
\bottomrule
\end{tabular}
\end{table}

A mortalidade hospitalar associada a esses procedimentos é
estruturalmente baixíssima ao longo de todo o período. Apenas
alguns casos isolados de óbito intra-hospitalar são registrados
ao longo de toda a série de dezenas de milhares de internações;
a taxa nacional anual oscila entre 0\,\% e 0{,}05\,\%. Isso
ratifica a propriedade clínica de \textit{baixo risco} explorada
na Seção~\ref{sec:teoria}: as diferenças de acesso \textit{não
podem} ser explicadas por seleção de risco do paciente.

\subsection{Desigualdade territorial: o gradiente
descritivo}\label{sub:desigualdade}

A Tabela~\ref{tab:uf-100k} apresenta a taxa por 100\,mil
habitantes em 2025, ordenada do maior ao menor.

\begin{table}[H]
\centering
\caption{Acesso à cirurgia vaginal de IU em 2025: AIH por
100\,mil habitantes, por UF. População IBGE/SIDRA tabela 6579.}
\label{tab:uf-100k}
\small
\begin{tabular}{lrrr}
\toprule
\textbf{UF} & \textbf{AIH 2025} & \textbf{Pop. (M)} & \textbf{AIH/100k} \\
\midrule
SC          & 1\,204 &  8{,}19 & 14{,}71 \\
PR          & 1\,381 & 11{,}89 & 11{,}61 \\
GO          &   621  &  7{,}42 &  8{,}37 \\
MS          &   198  &  2{,}92 &  6{,}77 \\
SP          & 2\,847 & 46{,}08 &  6{,}18 \\
ES          &   237  &  4{,}13 &  5{,}74 \\
RS          &   638  & 11{,}23 &  5{,}68 \\
TO          &    87  &  1{,}59 &  5{,}48 \\
MT          &   160  &  3{,}89 &  4{,}11 \\
MG          &   816  & 21{,}39 &  3{,}81 \\
AM          &   155  &  4{,}32 &  3{,}59 \\
MA          &   251  &  7{,}02 &  3{,}58 \\
RO          &    59  &  1{,}75 &  3{,}37 \\
PI          &   111  &  3{,}38 &  3{,}28 \\
BA          &   457  & 14{,}87 &  3{,}07 \\
CE          &   195  &  9{,}27 &  2{,}10 \\
DF          &    62  &  3{,}00 &  2{,}07 \\
RN          &    69  &  3{,}46 &  2{,}00 \\
AC          &    17  &  0{,}88 &  1{,}92 \\
RJ          &   317  & 17{,}22 &  1{,}84 \\
SE          &    39  &  2{,}30 &  1{,}70 \\
PB          &    67  &  4{,}16 &  1{,}61 \\
AP          &    12  &  0{,}81 &  1{,}49 \\
PE          &    70  &  9{,}56 &  0{,}73 \\
PA          &    63  &  8{,}71 &  0{,}72 \\
AL          &    10  &  3{,}22 &  0{,}31 \\
RR          &     1  &  0{,}74 &  0{,}14 \\
\bottomrule
\end{tabular}
\end{table}

A diferença entre o topo e a base é dramática: SC opera 14{,}71
cirurgias por 100\,mil habitantes em 2025, enquanto RR opera
0{,}14 \textemdash{} aproximadamente cem vezes menos. Não há
nenhum mecanismo demográfico ou epidemiológico plausível capaz
de explicar uma diferença de 100 vezes na demanda real por
cirurgia uroginecológica entre dois estados. A estrutura etária
da população feminina, a prevalência subjacente de IU e a
indicação cirúrgica são, em primeira aproximação, similares.
\textit{A diferença observada é desigualdade de oferta.}

A região Sul lidera consistentemente. SC e PR são, para todos os
fins práticos, \textit{outliers de cima}: superam não apenas
suas vizinhas geográficas mas também SP, que tem maior porte
hospitalar absoluto. Essa liderança é estrutural \textemdash{}
SC e PR ocupam topo do \textit{ranking} desde pelo menos 2015.
Estados periféricos do Norte (RR, AP), pequenos do Nordeste
(AL, SE) e \textit{regredidos} pós-pandemia (PE, RJ) compõem
a base. RJ e PE são particularmente preocupantes porque são
estados populosos e com tradição de formação médica: a
desigualdade ali não é trivial de explicar por insuficiência
demográfica.

\subsection{Pobreza estrutural por UF (cobertura
PBF)}\label{sub:pobreza}

A cobertura do Bolsa Família (\%~da população estadual
beneficiária) em 2024\textendash 2025 (Tabela~\ref{tab:cruzamento},
coluna \%\,Pop.\,em PBF) apresenta um gradiente bem conhecido,
mas que vale documentar lado a lado com o de acesso à cirurgia.
SC, RS, PR e SP operam com 3 a 6\,\% da população em PBF.
PA, BA, PE, MA, AL, PI operam com 17 a 19\,\%. Essa diferença
de cobertura captura, em parte, o gradiente de pobreza
estrutural brasileiro: estados do Sul e Sudeste com renda
mediana mais alta têm menor demanda formal pelo programa;
estados do Norte e Nordeste com renda mediana mais baixa têm
maior cobertura. O desvio padrão da cobertura PBF entre as 27
UFs é de aproximadamente 5{,}6 pontos percentuais; a média é de
12{,}3\,\%.

\subsection{Fluxo fiscal compensatório (emendas
parlamentares per capita)}\label{sub:emendas}

A execução de emendas parlamentares \textit{per capita}
(Tabela~\ref{tab:cruzamento}, coluna Emendas\,/\,hab.) opera em
\textit{outra direção} relativamente ao gradiente de pobreza.
Estados pequenos e\,/\,ou periféricos (AP \rs 608, RR \rs 442,
AC \rs 401) recebem volumes per capita várias vezes maiores que
estados grandes e ricos (SP \rs 44, RJ \rs 59, DF \rs 46). Isso
não é casual: o sistema de emendas tem um componente
\textit{compensatório implícito} \textemdash{} o cálculo
político da distribuição leva em conta tamanho da bancada
estadual relativamente à população, o que beneficia, em termos
\textit{per capita}, estados pequenos. Nominalmente, isso
poderia funcionar como um mecanismo redistributivo: estados
mais pobres recebem mais.

A pergunta deste artigo, contudo, não é \textit{normativa} (se
o sistema é justo) mas \textit{empírica} (se o fluxo extra
\textit{produz} oferta). É o que a Seção~\ref{sub:cruzamento}
investiga.

\subsection{O cruzamento: acesso, pobreza e fluxo
fiscal}\label{sub:cruzamento}

A Tabela~\ref{tab:cruzamento} consolida os três indicadores por
UF, em médias 2024\textendash 2025. Estados ordenados de cima
para baixo por acesso (AIH/100k).

\begin{table}[H]
\centering
\caption{Cruzamento entre acesso à cirurgia uroginecológica,
pobreza e fluxo fiscal por UF. Médias anuais 2024\textendash 2025.
Fontes: SIH-AIH-RD (UroPro), CGU\,/\,Portal da Transparência
(PBF e Emendas), IBGE/SIDRA, BCB (IPCA).}
\label{tab:cruzamento}
\small
\begin{tabular}{lrrr}
\toprule
\textbf{UF} & \textbf{AIH/100k} & \textbf{\%\,Pop.\,em PBF} & \textbf{Emendas\,/\,hab. (R\$\,2021)} \\
\midrule
SC &  12{,}33 &  3{,}4\,\% &   97 \\
PR &  10{,}68 &  5{,}9\,\% &   75 \\
GO &   7{,}69 &  7{,}6\,\% &  106 \\
MS &   7{,}25 &  7{,}9\,\% &  148 \\
SP &   6{,}04 &  6{,}2\,\% &   44 \\
TO &   4{,}98 & 11{,}1\,\% &  261 \\
RS &   4{,}96 &  6{,}3\,\% &   93 \\
MT &   4{,}37 &  7{,}3\,\% &  128 \\
AC &   4{,}30 & 16{,}5\,\% &  401 \\
ES &   4{,}29 &  8{,}5\,\% &  119 \\
MA &   4{,}12 & 19{,}0\,\% &  135 \\
RO &   3{,}74 &  8{,}8\,\% &  226 \\
MG &   3{,}72 &  8{,}3\,\% &   78 \\
AM &   3{,}39 & 16{,}4\,\% &  154 \\
BA &   2{,}73 & 18{,}0\,\% &   91 \\
PI &   2{,}73 & 18{,}9\,\% &  223 \\
DF &   2{,}49 &  6{,}5\,\% &   46 \\
CE &   2{,}23 & 16{,}9\,\% &  111 \\
RN &   2{,}04 & 15{,}6\,\% &  158 \\
RJ &   1{,}71 & 10{,}1\,\% &   59 \\
PB &   1{,}44 & 17{,}4\,\% &  157 \\
SE &   1{,}39 & 17{,}9\,\% &  216 \\
PE &   0{,}82 & 18{,}0\,\% &  102 \\
AP &   0{,}74 & 16{,}6\,\% &  608 \\
PA &   0{,}66 & 16{,}9\,\% &  108 \\
RR &   0{,}54 & 12{,}8\,\% &  442 \\
AL &   0{,}31 & 18{,}1\,\% &  181 \\
\bottomrule
\end{tabular}
\end{table}

A leitura cruzada da tabela rende três achados-chave.

\textbf{Achado 1 \textemdash{} Pobreza e acesso são fortemente
correlacionados em sentido inverso.} Os cinco estados de maior
acesso (SC, PR, GO, MS, SP) operam com cobertura PBF entre
3{,}4 e 7{,}9\,\%. Os cinco de menor acesso (AL, RR, PA, AP,
PE) operam com cobertura entre 12{,}8 e 18{,}1\,\%. A correlação
de Pearson entre as duas séries é de aproximadamente
\textendash 0{,}68 sobre as 27 UFs. Em outras palavras: cada
ponto percentual a mais de cobertura PBF em uma UF está
associado, no agregado nacional, a uma redução substancial na
taxa de cirurgias uroginecológicas por 100\,mil habitantes. O
sentido da correlação é o esperado pela teoria
(pobreza\,$\rightarrow$\,menos serviços especializados); a
\textit{magnitude} da correlação é alta para qualquer padrão
de dado público brasileiro, e isso é o que merece ênfase.

\textbf{Achado 2 \textemdash{} O fluxo fiscal compensatório não
compensa.} A correlação entre emendas \textit{per capita} e
acesso à cirurgia é de aproximadamente \textendash 0{,}45
\textemdash{} \textit{também negativa}. Estados que recebem
mais emendas \textit{per capita} são, em média, os de menor
acesso. AP recebe \rs 608 \textit{por habitante por ano} em
emendas (em moeda constante 2021) e operou apenas 0{,}74
cirurgias\,/\,100k em 2025. SP recebe \rs 44 e opera 6{,}04.
A leitura mais conservadora desse padrão é que o fluxo
político-fiscal das emendas tem componente compensatório
\textit{de gasto}, mas não \textit{produz oferta hospitalar
especializada visível}. As emendas são, em média, alocadas em
obras locais, custeio genérico e capilarização infraestrutural;
a montagem de uma equipe cirúrgica especializada em
uroginecologia exige investimentos qualificados, treinamento
contínuo e estrutura hospitalar terciária que volume nominal
\textit{per capita} de gasto não compra automaticamente. O dado
é coerente com a tese geral de que \textit{capacidade}
hospitalar especializada é uma função de longo prazo, não
proporcional ao orçamento corrente.

\textbf{Achado 3 \textemdash{} Existem três UFs que ``driblam''
o gradiente.} Tocantins, Acre e Maranhão operam acesso
sensivelmente mais alto do que esperaríamos pelo seu nível de
pobreza estrutural. TO tem 11\,\% da população em PBF e atinge
4{,}98 AIH/100k (acima da média nacional de aproximadamente
4{,}1). AC tem 16{,}5\,\% e atinge 4{,}30; MA tem 19\,\% e
atinge 4{,}12. Em contraste, BA (18\,\% PBF, 2{,}73), PE (18\,\%
PBF, 0{,}82) e AL (18\,\% PBF, 0{,}31) operam acesso muito
abaixo. A diferença entre MA e AL, ambos com cobertura PBF
similar, é de aproximadamente 13$\times$ no acesso à cirurgia.

A existência dessas \textit{exceções} é metodologicamente
importante. Mostra que o gradiente pobreza\,$\rightarrow$\,acesso
\textit{não é determinístico}. Configurações organizacionais
locais \textemdash{} como a rede pública estadual está
estruturada, quais hospitais foram credenciados, qual é a
política estadual de capacitação cirúrgica, qual o histórico de
relacionamento com universidades médicas \textemdash{} podem
\textit{driblar} o efeito da pobreza estrutural. Identificar e
caracterizar essas exceções é uma das agendas mais promissoras
de política pública brasileira em saúde, e está fora do escopo
descritivo deste paper.

\subsection{O choque pandêmico como teste natural}\label{sub:covid}

A pandemia de COVID-19 funcionou, no caso aqui examinado, como
um \textit{teste de carga} sobre o sistema federativo. A queda
de volume nacional foi homogênea entre UFs (em torno de
\textendash 50\textendash 70\,\% em quase todos os estados em
2020\textendash 2021); a \textit{recuperação}, contudo, foi
fortemente heterogênea. SC e PR, ambos no topo do
\textit{ranking} de acesso pré-pandêmico, expandiram de forma
expressiva (+138\,\% e +84\,\% no volume médio
2022\textendash 2025 \textit{vs.} 2015\textendash 2019); RJ
(\textendash 3\,\%), PE (\textendash 35\,\%), DF (\textendash 23\,\%)
e AP (\textendash 81\,\%) sequer voltaram ao patamar
pré-pandêmico.

Sob a lente do argumento principal deste paper, o choque
pandêmico é particularmente revelador: ele \textit{amplificou
mas não criou} a desigualdade preexistente. UFs de oferta forte
absorveram represa cirúrgica e seguiram crescendo; UFs de oferta
fraca regrediram ou não recuperaram. A pandemia, em outras
palavras, expôs e quantificou diferenças de robustez sistêmica
que já existiam, e tornou ainda mais visível o gradiente
pobreza\,$\rightarrow$\,acesso documentado nas seções anteriores.

% =====================================================================
\section{Discussão}\label{sec:discussao}
% =====================================================================

\subsection{O federalismo brasileiro e a tensão
universalismo\,/\,distribuição}

O achado central deste paper \textemdash{} de que o acesso à
cirurgia uroginecológica acompanha de perto o gradiente de
pobreza estadual e de que as emendas parlamentares
\textit{per capita} não compensam essa desigualdade
\textemdash{} é coerente com uma leitura clássica do federalismo
brasileiro: o desenho institucional de 1988 produz
\textit{universalismo formal} (todo cidadão tem direito ao SUS)
mas mantém uma \textit{distribuição material altamente desigual}
da capacidade instalada. A oferta de serviços hospitalares
especializados é função de fatores estruturais de longo prazo
\textemdash{} densidade da rede, formação médica, fluxo de
referência, governança estadual da rede \textemdash{} cuja
correlação com pobreza estrutural é forte e estável. Mecanismos
fiscais de curto e médio prazo, como as emendas parlamentares,
podem alterar volumes nominais de gasto sem modificar a
distribuição de oferta clínica especializada.

Esse não é um argumento contra a existência de emendas
parlamentares como mecanismo, mas contra a \textit{inferência
direta} \textemdash{} comum no debate público e legislativo
\textemdash{} de que o volume \textit{per capita} de emendas
recebido por uma UF se traduz em mais e melhores serviços
públicos. O dado mostra que, ao menos para esse procedimento e
nessa janela temporal, a tradução não acontece.

\subsection{Por que as emendas não viram oferta cirúrgica
especializada?}

Cinco hipóteses, não-mutuamente exclusivas, articulam por que o
fluxo de emendas \textit{per capita} não compensa o gradiente
de oferta cirúrgica:

\textbf{(i) Inadequação setorial.} As emendas, em média, não
são alocadas para construção e manutenção de redes hospitalares
de média e alta complexidade; tendem a financiar obras
municipais, custeio genérico de unidades básicas e ações de
visibilidade política (eventos, equipamentos pontuais).

\textbf{(ii) Inadequação temporal.} Mesmo emendas alocadas a
saúde tendem a financiar gastos de \textit{curto prazo}; a
montagem de uma equipe cirúrgica especializada exige
investimentos plurianuais que vão além do horizonte de mandato
parlamentar.

\textbf{(iii) Capacidade absortiva local.} Mesmo onde há
recurso disponível, há UFs cuja rede pública estadual não tem
capacidade institucional de transformar gasto em capacidade
clínica especializada (déficit de gestão, fragmentação
hospitalar, ausência de programa próprio de capacitação).

\textbf{(iv) Capilaridade vs. concentração.} Procedimentos
uroginecológicos exigem hospitais com porte e
multidisciplinaridade. Emendas \textit{per capita} altas em UFs
pequenas frequentemente se traduzem em \textit{capilarização}
da rede (mais postos, mais unidades) sem aumentar o número de
hospitais com porte suficiente para uroginecologia.

\textbf{(v) Capacidade de exportação para fora do estado.}
UFs de oferta baixa frequentemente operam Tratamento Fora do
Domicílio (TFD): pacientes são deslocados para grandes centros
de outros estados. Isso significa que a oferta nominal local
não captura todo o atendimento dos cidadãos da UF, mas também
\textit{transfere recurso fiscal a outros estados} sem
desenvolver capacidade local.

A caracterização empírica de qual dessas hipóteses domina em
cada UF é uma agenda de pesquisa relevante e está fora do
escopo deste paper. O importante para o argumento aqui é que
\textit{nenhuma} delas é plausivelmente corrigida por aumento
do volume nominal de emendas, e portanto a expectativa
político-discursiva de que ``mais emendas'' equivalha a
``melhor oferta'' é, no mínimo, otimista.

\subsection{O que TO, AC e MA fazem de diferente?}

A existência de três UFs que driblam o gradiente é, talvez, o
achado mais promissor do ponto de vista de política pública.
Esses três estados têm cobertura PBF entre 11 e 19\,\% (ou
seja, são estruturalmente pobres) mas operam taxas de acesso
\textit{na média ou acima da média nacional}. A questão
relevante é: o que essas configurações estaduais têm em comum
que outras UFs igualmente pobres (AL, PE, BA) não têm?

Sem dados primários sobre organização da rede pública
estadual, levantamos apenas hipóteses qualitativas, todas
testáveis em pesquisa subsequente. \textbf{(a)} Tocantins,
desde sua criação em 1988, fez investimentos significativos em
rede hospitalar pública estadual concentrada em Palmas, com
referência consolidada para os procedimentos eletivos da região;
isso pode ter criado capacidade local desproporcional ao
tamanho e à riqueza estadual. \textbf{(b)} Acre tem uma rede
hospitalar pública pequena mas tecnicamente forte, com
relacionamento histórico com universidades de outros estados,
o que pode ter sustentado capacidade especializada apesar da
pobreza estrutural. \textbf{(c)} Maranhão investiu, na década
2010\textendash 2020, em programa estadual de hospitais
referenciados, com resultado mensurável em vários indicadores
de saúde. Essas três configurações são, do ponto de vista de
política pública, \textit{exemplares} de como contornar o
gradiente; deveriam ser objeto de estudos de caso aprofundados,
para tipologizar o que é replicável e o que é específico.

\subsection{Considerações sobre o defeito do pipeline}

Vale tratar explicitamente um ponto que fica implícito ao longo
do paper. As tabelas e correlações apresentadas refletem dados
\textit{posteriores} à correção descrita na
Seção~\ref{sub:bug}: o pipeline silver da plataforma Mirante
operava, até abril de 2026, com um filtro defeituoso que
ocultava silenciosamente 14 das 27 UFs. Os mapas e gráficos da
plataforma mostravam, antes da correção e sem qualquer aviso,
uma versão \textit{ainda mais desigual} da realidade do que
esta corrigida.

Essa observação tem dupla relevância. \textbf{Primeiro}, sublinha
um ponto cético sobre a credibilidade do dado público: a
publicação do dado bruto pelo DATASUS é necessária mas \textit{não
suficiente}; a engenharia de dados \textit{a jusante} é parte
integral da infraestrutura, e seus defeitos podem inverter
conclusões. \textbf{Segundo}, demonstra a vantagem de pipelines
abertos e auditáveis: o defeito foi detectado por inspeção visual
de um mapa, diagnosticado por consulta SQL direta ao Delta da
camada \textit{silver}, corrigido por um \textit{commit} de uma
linha removida e um comentário substituído, e teve seu impacto
quantificado pela regeneração das camadas \textit{silver} e
\textit{gold}. Toda essa cadeia foi pública, reprodutível e
versionada. Em arquiteturas opacas (planilhas adhoc, ETLs em
ferramentas \textit{black-box}, dashboards proprietários), esse
mesmo bug poderia ter persistido por anos sem ser detectado.

\subsection{Limitações de inferência}

(i)~Correlação não é causalidade. Que cobertura PBF correlacione
\textendash 0{,}68 com acesso à cirurgia não significa que
\textit{aumentar} PBF \textit{reduziria} cirurgias \textemdash{}
é apenas que ambas refletem um substrato comum (pobreza
estrutural). (ii)~Emendas alocadas a saúde \textit{específicamente}
não foram destacadas neste paper; é possível que o subconjunto
de emendas saúde tenha relação diferente com oferta cirúrgica.
(iii)~A análise é estática (médias 2024\textendash 2025). Uma
análise dinâmica (mudanças no acesso após mudanças no fluxo de
emendas, ou após mudanças na cobertura PBF) seria mais robusta
para inferência causal. (iv)~A IU cirúrgica é apenas \textit{um}
indicador; a generalização da conclusão para outros
procedimentos eletivos exige replicação.

% =====================================================================
\section{Considerações Finais}\label{sec:conclusao}
% =====================================================================

A cirurgia uroginecológica para tratamento da incontinência
urinária é, no Brasil contemporâneo, um procedimento clínico
\textit{tecnicamente maduro} \textemdash{} a permanência
hospitalar média caiu 40\,\% na via vaginal entre 2008 e 2025,
a mortalidade é estruturalmente baixa, a difusão dos
\textit{slings} sub-uretrais alinha o país aos padrões
internacionais. Essa maturidade técnica, contudo, não se
traduziu em equidade de acesso. Em 2025, a taxa por 100\,mil
habitantes é cem vezes maior em Santa Catarina do que em
Roraima.

O cruzamento com indicadores estaduais de pobreza estrutural
(cobertura do Bolsa Família) e de fluxo fiscal compensatório
(emendas parlamentares \textit{per capita}) mostra que essa
desigualdade não é aleatória. Ela acompanha de perto o
gradiente de pobreza ($\rho \approx -0{,}68$) e \textit{não é}
corrigida pelo mecanismo de emendas \textemdash{} pelo
contrário, emendas \textit{per capita} maiores estão
correlacionadas com \textit{menor} acesso ($\rho \approx -0{,}45$).
A imagem que emerge é a de um sistema federativo em que a
oferta de serviços hospitalares especializados é uma função
estrutural de longo prazo, e em que o fluxo político-fiscal de
curto prazo, embora volumoso em UFs periféricas, não consegue
produzir capacidade de oferta visível.

Há, porém, três UFs \textemdash{} TO, AC e MA \textemdash{} que
operam acima do esperado dado seu nível de pobreza, sugerindo
que existem configurações organizacionais locais capazes de
contornar o gradiente. Caracterizar essas configurações é,
provavelmente, a agenda de política pública mais promissora
sugerida por este trabalho.

Em termos metodológicos, o artigo registra (a)~a viabilidade de
análise \textit{cross-vertical} sobre dados públicos brasileiros
quando passados por uma arquitetura unificada de \textit{pipeline}
aberto, e (b)~a importância da auditabilidade da camada
\textit{silver} para a credibilidade do dado entregue ao público.
O defeito específico documentado (filtro
\texttt{== max(\_ingest\_ts)}) e sua correção em \textit{commit}
\texttt{fa869cf} são, em si, contribuição metodológica modesta
mas concreta: defeitos similares são plausivelmente comuns em
\textit{pipelines} de produção construídos em cima de Auto Loader,
e tornar o caso público pode ajudar outras equipes a evitá-lo.

Universalismo formal e desigualdade material continuam, em 2025,
coexistindo no SUS. A constatação não é nova; o que esperamos
contribuir aqui é uma quantificação direta, em um procedimento
clínico bem definido, de quão acentuada permanece a desigualdade,
e de como o instrumento institucional principal de
\textit{compensação fiscal explícita} \textemdash{} as emendas
parlamentares \textemdash{} não corrige o gradiente. Se a aspiração
universalista do SUS é levada a sério, mecanismos
\textit{adicionais} àqueles já em uso \textemdash{} qualificação
de redes estaduais, programas de capacitação cirúrgica
descentralizada, integração com formação médica regional
\textemdash{} parecem necessários.

% =====================================================================
\section*{Referências}\label{sec:ref}
\addcontentsline{toc}{section}{Referências}
% =====================================================================
\begin{singlespace}\small
\noindent\hangindent=1.25cm\hangafter=1
ARMBRUST, M.; GHODSI, A.; XIN, R.; ZAHARIA, M. Lakehouse: a new
generation of open platforms that unify data warehousing and
advanced analytics. In: \textbf{11th CIDR}, 2021.\par\smallskip

\noindent\hangindent=1.25cm\hangafter=1
BANCO CENTRAL DO BRASIL (BCB). \textbf{Sistema Gerenciador de
Séries Temporais}, série 433: IPCA. Disponível em:
\url{https://api.bcb.gov.br/dados/serie/bcdata.sgs.433/dados}.
Acesso em: abr. 2026.\par\smallskip

\noindent\hangindent=1.25cm\hangafter=1
BRASIL. Ministério da Saúde. \textbf{Sistema de Informações
Hospitalares do SUS (SIH-SUS)}. Microdados. Brasília: DATASUS.
Disponível em:
\url{https://datasus.saude.gov.br/transferencia-de-arquivos/}.
Acesso em: abr. 2026.\par\smallskip

\noindent\hangindent=1.25cm\hangafter=1
BRASIL. Ministério da Saúde. \textbf{Tabela de Procedimentos,
Medicamentos e OPM do SUS \textemdash{} SIGTAP}. Brasília:
DATASUS, 2024.\par\smallskip

\noindent\hangindent=1.25cm\hangafter=1
CHALHOUB, L. Emendas Parlamentares no Orçamento Federal
Brasileiro (2015\textendash 2025): distribuição espacial,
execução orçamentária e efeitos das mudanças institucionais
recentes. \textbf{Mirante dos Dados \textemdash{} Working Paper
n. 1}, abr. 2026a.\par\smallskip

\noindent\hangindent=1.25cm\hangafter=1
CHALHOUB, L. Programa Bolsa Família, Auxílio Brasil e Novo
Bolsa Família (2013\textendash 2025): transformações
institucionais, expansão da cobertura e desigualdade territorial.
\textbf{Mirante dos Dados \textemdash{} Working Paper n. 2},
abr. 2026b.\par\smallskip

\noindent\hangindent=1.25cm\hangafter=1
CHAPPLE, C. R.; CRUZ, F.; DESCHAMPS, C.; HAYLEN, B. T.
\textit{et al.} Consensus statement of the European Association
of Urology on the surgical treatment of urinary incontinence in
women. \textbf{European Urology}, v. 78, n. 5,
p. 643\textendash 656, 2020.\par\smallskip

\noindent\hangindent=1.25cm\hangafter=1
CHALHOUB, L. Cirurgia uroginecológica no Sistema Único de Saúde,
2008\textendash 2025: ganhos silenciosos de eficiência,
desigualdade territorial persistente, choque pandêmico e represa
cirúrgica. \textbf{Mirante dos Dados \textemdash{} Working
Paper n. 5}, abr. 2026f.\par\smallskip

\noindent\hangindent=1.25cm\hangafter=1
FORD, A. A.; ROGERSON, L.; CODY, J. D.; OGAH, J. Mid-urethral
sling operations for stress urinary incontinence in women.
\textbf{Cochrane Database of Systematic Reviews}, n. 7,
CD006375, 2017.\par\smallskip

\noindent\hangindent=1.25cm\hangafter=1
HAYLEN, B. T.; DE RIDDER, D.; FREEMAN, R. M. \textit{et al.}
An International Urogynecological Association
(IUGA)\,/\,International Continence Society (ICS) joint report
on the terminology for female pelvic floor dysfunction.
\textbf{Neurourology and Urodynamics}, v. 29, n. 1,
p. 4\textendash 20, 2010.\par\smallskip

\noindent\hangindent=1.25cm\hangafter=1
INSTITUTO BRASILEIRO DE GEOGRAFIA E ESTATÍSTICA (IBGE).
\textbf{SIDRA \textemdash{} Sistema IBGE de Recuperação Automática}.
Tabela 6579: Estimativas anuais da população residente.
Disponível em: \url{https://sidra.ibge.gov.br/tabela/6579}.
Acesso em: abr. 2026.\par\smallskip

\noindent\hangindent=1.25cm\hangafter=1
NUNES, B. P.; FLORES, T. R.; MIELKE, G. I. \textit{et al.} The
COVID-19 pandemic and elective surgeries in the Brazilian
Unified Health System (SUS). \textbf{Public Health}, v. 211,
p. 109\textendash 115, 2022.\par\smallskip

\noindent\hangindent=1.25cm\hangafter=1
TATIELI DA SILVA. \textbf{Análise Descritiva da Cirurgia para
Tratamento da Incontinência Urinária no SUS, 2015\textendash 2020}.
Trabalho de Conclusão de Curso (Especialização em Enfermagem),
2022.\par\smallskip

\noindent\hangindent=1.25cm\hangafter=1
WILKINSON, M. D. \textit{et al.} The FAIR Guiding Principles
for scientific data management and stewardship.
\textbf{Scientific Data}, v. 3, 2016.
DOI: 10.1038/sdata.2016.18.\par\smallskip
\end{singlespace}

\vspace{1cm}

\noindent\textbf{Citar como:}\\
CHALHOUB, L. Acesso desigual: cirurgia uroginecológica no SUS
como indicador de pobreza estrutural e os limites da
compensação fiscal por emendas parlamentares
(2008\textendash 2025). \textit{Mirante dos Dados \textemdash{}
Working Paper n. 3}, v. 2.0, abr. 2026. Disponível em:
\url{https://leonardochalhoub.github.io/mirante-dos-dados-br/}.

\vspace{0.5cm}

\noindent\textbf{Reconhecimento:} este Working Paper tem como
base o trabalho de especialização em Enfermagem de TATIELI
(2022). A presente versão reorienta o argumento em torno do
cruzamento
\textit{cross-vertical} entre cirurgia uroginecológica, pobreza
estrutural (Bolsa Família, Mirante WP n. 2) e fluxo fiscal
político (Emendas, Mirante WP n. 1), e incorpora a correção da
camada \textit{silver} (\textit{commit} \texttt{fa869cf}, abril
de 2026).

\vspace{0.5cm}

\noindent\textbf{Licença:} dados consolidados (camada
\textit{gold}) e código do \textit{pipeline} distribuídos sob
licença MIT. Texto deste artigo distribuído sob licença Creative
Commons Atribuição 4.0 Internacional (CC BY 4.0).

\end{document}
